\documentclass[11pt]{article}

\usepackage[utf8]{inputenc}
\usepackage[T1]{fontenc}
\usepackage{lmodern}
% \usepackage[margin=1in]{geometry}
\usepackage{hyperref}
\usepackage{parskip}
\usepackage{amsmath}
\usepackage{amssymb}
\usepackage{float}

\newcommand{\dnn}{\boldsymbol{\mathsf N}}
\newcommand{\dnnlayer}[1][i]{\dnn_{\boldsymbol{\hat\theta}^{(#1)}}}
\newcommand{\D}{\boldsymbol{\mathsf D}}
\makeatletter
\newcommand{\@varsup}[1]{\@ifnextchar\bgroup{\@varsup@arg{#1}}{#1}}
\newcommand{\@varsup@arg}[2]{#1^{(#2)}}
\newcommand{\residual}{\@varsup{\boldsymbol r}}
\newcommand{\x}{\@varsup{\boldsymbol x}}
\newcommand{\estimate}{\@varsup{\boldsymbol{\hat x}}}
\newcommand{\dirty}{\boldsymbol x_{\text{d}}}
\newcommand{\true}{\boldsymbol{x^{*}}}
\newcommand{\noise}{\boldsymbol{b}}
\makeatother

\title{Finding failure modes of R2D2 and CLEAN}
\author{Thomas Harper}
\date{\today}

\begin{document}

\maketitle
\newpage
\tableofcontents
\newpage

\section{Overview}

We search for failure modes by treating the observing setup as a parameter vector
\begin{equation}
\boldsymbol\theta =
(\rho_{\text{DR}}, t_{\text{obs}}, \nu_0, \Delta\nu, C, \ldots),
\end{equation}
\begin{itemize}
    \item $\boldsymbol\theta$ is the observing setup searched by nested sampling
    \item $\rho_{\text{DR}}$ is the target dynamic range
    \item $t_{\text{obs}}$ is the observation duration
    \item $\nu_0$ is the start frequency
    \item $\Delta\nu$ is the channel width
    \item $C$ is the telescope configuration
\end{itemize}
then we simulate data $\boldsymbol y(\boldsymbol\theta)$ using the MeqTrees software and attempt to reconstruct it with R2D2 and CLEAN.
For an algorithm $A$, nested sampling targets a failure score
\begin{equation}
F_A(\boldsymbol\theta)
=
\mathcal M\left(\estimate_A(\boldsymbol y(\boldsymbol\theta)), \true(\boldsymbol\theta)\right),
\end{equation}
\begin{itemize}
    \item $A$ is the reconstruction algorithm, either R2D2 or CLEAN
    \item $F_A(\boldsymbol\theta)$ is the failure score for algorithm $A$ at setup $\boldsymbol\theta$
    \item $\boldsymbol y(\boldsymbol\theta)$ is the simulated visibility data
    \item $\estimate_A$ is the reconstruction produced by algorithm $A$
    \item $\true(\boldsymbol\theta)$ is the ground-truth sky (for the moment simply a point source)
    \item $\mathcal M$ is the chosen metric or composite failure score
\end{itemize}
where larger $F_A$ means worse reconstruction, higher residuals, or greater computational cost.
The high-scoring samples identify regions of observing conditions where R2D2, CLEAN, or both are likely to fail.

\section{Data Model}
\label{sec:data-model}

\begin{equation}
\boldsymbol y = \boldsymbol\Phi \true + \boldsymbol n
\end{equation}
\begin{itemize}
    \item $\boldsymbol y\in \mathbb C^M$ is the measured visibilities
    \item $\boldsymbol\Phi: \mathbb R^N \to \mathbb C^M$  is the measurement operator
    \item $\true\in\mathbb R_+^N$ is the true sky image
    \item $\boldsymbol n\in\mathbb C^M$ is Gaussian noise with s.d. $\tau>0$ and mean $0$
\end{itemize}

\begin{equation}
\dirty = \kappa\operatorname{Re}\{\boldsymbol\Phi^\dagger \boldsymbol y\} = \D\true + \noise
\end{equation}

\begin{itemize}
    \item $\dirty$ is the dirty image including noise
    \item $\kappa=\max(\operatorname{Re}{\boldsymbol\Phi^\dagger\boldsymbol\Phi\boldsymbol\delta})^{-1}>0$ is the normalisation factor ensuring the PSF ($\boldsymbol{h}=\kappa \operatorname{Re}\left\{\boldsymbol{\Phi}^{\dagger} \boldsymbol{\Phi} \boldsymbol{\delta}\right\} \in \mathbb{R}^N$) has peak value 1, where $\boldsymbol{\mathbf\delta} \in \mathbb R^N$ a discrete Dirac delta centred in the image (1 at the centre, 0 elsewhere)
    \item $\boldsymbol\Phi^\dagger$ is the adjoint of the measurement operator
    \item $\D:\mathbb R^N\to\mathbb R^N,\quad \D \triangleq \kappa\operatorname{Re}\{\boldsymbol{\Phi}^{\dagger}\boldsymbol{\Phi}\}$ is a normalised map for the true image to the noise free dirty image, encoding the non-uniform Fourier sampling, gridding, de-gridding and Fourier transform operations
    \item $\D\true$ is the noise-free dirty image
    \item $\noise = \kappa\operatorname{Re}\{\boldsymbol\Phi^\dagger\boldsymbol n\}\in\mathbf R^N$ is the normalised back projected noise
\end{itemize}

\section{Reconstruction Algorithms}

\subsection{R2D2}

\begin{gather}
    \residual{i} = \dirty - \D\x{i}\\
    \x{i} = \x{i-1} + \dnnlayer(\residual{i-1}, \x{i-1})
\end{gather}

\begin{itemize}
    \item $\residual{i}$ is a residual dirty image. It is the difference between the original noisy dirty image $\dirty$ and the one derived from the $i^\text{th}$ estimation of the night sky $\D\x{i}$
    \item $\x{i}$ is the $i^\text{th}$ estimate of the night sky produced by R2D2
    \item $\dnnlayer$ is the $i^\text{th}$ deep neural network layer in R2D2
    \item $\boldsymbol{\hat\theta}^{(i)}\in\mathbb R^Q$ are the learned parameters for a given DNN layer.
\end{itemize}

\begin{equation}
\estimate \triangleq \x{I} = \sum_{i=1}^I\dnnlayer(\residual{i-1}, \x{i-1})
\end{equation}

\begin{itemize}
    \item $\estimate$ is R2D2s final reconstruction
    \item $I$ is the total number of DNN layers
\end{itemize}

\subsection{Högbom CLEAN}

\begin{gather}
    \residual{i} = \dirty - \D\x{i}\\
    \x{i} = \x{i-1} + \gamma\,\residual{i-1}_{p^{(i)}}\,\boldsymbol{\delta}_{p^{(i)}}
\end{gather}

\begin{itemize}
    \item $\residual{i}$ is a residual dirty image. It is the difference between the original noisy dirty image $\dirty$ and the one derived from the $i^\text{th}$ CLEAN component model $\D\x{i}$
    \item $\x{i}$ is the $i^\text{th}$ CLEAN component model of the sky image, initialised at $\x{0}=\boldsymbol 0$
    \item $p^{(i)}=\arg\max_{n}\lvert\residual{i-1}_{n}\rvert$ is the pixel of the peak residual
    \item $\boldsymbol{\delta}_{p^{(i)}}$ is a discrete Dirac delta at pixel $p^{(i)}$
    \item $\gamma\in(0,1]$ is the loop gain
\end{itemize}

\begin{equation}
\x{I} = \sum_{i=1}^{I}\gamma\,\residual{i-1}_{p^{(i)}}\,\boldsymbol{\delta}_{p^{(i)}}
\end{equation}

\begin{equation}
\estimate \triangleq \boldsymbol g \ast \x{I} + \residual{I}
\end{equation}

\begin{itemize}
    \item $\x{I}$ is the CLEAN component model after $I$ iterations
    \item $\boldsymbol g$ is the restoring beam, typically a Gaussian fit to the main lobe of the PSF
    \item $\estimate$ is the restored CLEAN image
    \item $I$ is the number of CLEAN iterations
\end{itemize}

\section{Evaluation Metrics}

\subsection{Reconstruction Fidelity}

\subsubsection{Signal-to-Noise Ratio}

\begin{equation}
    \mathrm{SNR}(\estimate, \true)
    =
    20\log_{10}
    \left(
    \frac{\|\true\|_2}
    {\|\true - \estimate\|_2}
    \right)
\end{equation}

\subsubsection{Logarithmic Signal-to-Noise Ratio}

\paragraph{R2D2 definition}

\begin{gather}
    \operatorname{rlog}(\x)=x_{\max} \log _a\left(\frac{a}{x_{\max}} \x+\mathbf{1}\right)\\
    \operatorname{logSNR}(\estimate, \true)=\operatorname{SNR}(\operatorname{rlog}(\estimate), \operatorname{rlog}(\true))
\end{gather}

\begin{itemize}
    \item $a$ is a parameter set to $\rho_\text{DR}$, the dynamic range of the ground truth image
    \item $x_\text{max}$ is the peak pixel value of the image $\x$
    \item $\mathbf 1\in\mathbb R^N$ is a vector of values equal to 1
\end{itemize}

\paragraph{Standard definition}

\subsubsection{Dynamic Range}

\subsubsection{Normalised Mean Squared Error}

\begin{equation}
    \operatorname{NMSE}(\estimate, \true)
    =
    \frac{\|\estimate - \true\|_2^2}
    {\|\true\|_2^2}
\end{equation}

\begin{itemize}
    \item $\estimate$ is the reconstructed image
    \item $\true$ is the ground-truth sky image
    \item $\|\estimate - \true\|_2^2$ is the squared reconstruction error
    \item $\|\true\|_2^2$ normalises the error by the ground-truth image energy
\end{itemize}

Lower NMSE indicates better reconstruction fidelity.

\subsection{Image-Domain Data Fidelity}

\subsubsection{Normalised Residual}

$$
\bar\sigma_{\text{res.}}(\hat{\residual},\dirty)=\frac{\|\hat{\residual}\|_2}{\|\dirty\|_2}
$$

\begin{itemize}
\item $\hat\residual$ - final dirty image residual
\item $\dirty$ - true dirty image
\end{itemize}

As $\hat\x\to\true$, $\hat\residual\to0$, thus lower is better.

Essentially this says, normalised by the true dirty image $\dirty$, how much is left in the dirty image residual $\hat\residual$.
Ideally nothing is left because applying the dirty-image operator $\D$ to the predicted image $\estimate$ reproduces the dirty image $\dirty$, so $\D\estimate\approx\dirty$ and thus $\hat\residual=\dirty-\D\estimate\approx0$.

In short: how closely does each model reproduce the dirty image?

\subsection{Scientific Fidelity}

\subsubsection{Flux Recovery}

\subsubsection{Astrometric Error}

\subsubsection{Source Detection and Completeness}

\section{Search Parameter Space}

\subsection{Sky and Image Parameters}

\subsubsection{Dynamic Range}

In the nested-sampling search, dynamic range is a simulation parameter rather than a metric computed from the output image.
MeqTrees predicts visibilities for a fixed point-source sky, then the simulator adds complex Gaussian thermal noise with
\begin{equation}
    \sigma_{\text{vis}}
    =
    \frac{S_{\text{src}}}{\rho_{\text{DR}}}.
\end{equation}

\begin{itemize}
    \item $\rho_{\text{DR}}$ is the requested simulation dynamic range
    \item $S_{\text{src}}$ is the true point-source flux, fixed to $1~\mathrm{Jy}$ in the current runs
    \item $\sigma_{\text{vis}}$ is the complex visibility noise standard deviation
\end{itemize}

Thus lower $\rho_{\text{DR}}$ produces noisier visibilities before R2D2 or CLEAN sees the data.

\subsubsection{Source Morphology}

\subsubsection{Source Position}

\subsection{Observation and Fourier-Sampling Parameters}

\subsubsection{VLA Configuration}

\subsubsection{Declination}

\subsubsection{Observation Duration}

VLA dynamic-scheduling timescales, EVLA Memo~232 \cite{faes2024evla232}:
\begin{itemize}
    \item Scan length: 0--10~min (Fig.~7)
    \item Single tracking (same target, consecutive scans): 0--10~min (Fig.~11)
    \item Time on source per scheduling block (even if interrupted): 0--20~min (Fig.~17)
\end{itemize}

For regular (standard observing) scans, no less than 20~seconds on source after slewing (Static Report Tab) \cite{nrao-vla-obsguide}.
Suggested search range: 20~s--20~min.

\subsubsection{Observation Start Time and Hour-Angle Coverage}

\subsubsection{Number of Frequency Channels}

\subsubsection{Frequency Range and Bandwidth}

\begin{table}[H]
\centering
\begin{tabular}{ll}
4  & 54--86~MHz \\
P  & 224--480~MHz \\
L  & 1--2~GHz \\
S  & 2--4~GHz \\
C  & 4--8~GHz \\
X  & 8--12~GHz \\
Ku & 12--18~GHz \\
K  & 18--26.5~GHz \\
Ka & 26.5--40~GHz \\
Q  & 40--50~GHz \\
\end{tabular}
\caption{VLA receiver bands. Recommended calibrator cycle times by configuration are in Table~3.1 \cite{nrao-vla-obsguide}.}
\label{tab:vla-bands}
\end{table}

\subsubsection{Flagging Fraction}

\subsubsection{Visibility Weighting}

\section{Proposed Search}

\subsection{Search Variables}

\subsection{Parameter Ranges}

\subsection{Objective Metrics}

\subsection{Failure Criteria}

\nocite{*}
\bibliographystyle{plain}
\bibliography{references}

\end{document}
